% Template created by Steven Wolf (wolfs15@ecu.edu) heavily based on the template created by
% Adrienne Traxler.
% Last modified: 4/16/19: Added title page for new format.  Made small edits to the document
% text to reflect changes in PERC proceedings formatting guidelines.
% 
% Changelog
% 4/16/19:  Added title page for new format.  Made small edits to the document text to reflect
% changes in PERC proceedings formatting guidelines.
% 3/14/17: add 2-column figure example, tweak text to align with Word sample file, add hyperref for links/cross-references
% 6/21/16: 	Fix old PACS option and superscript affiliations.
% 5/5/16: 	Update sample table to use the REVTEX ruledtabular environment for auto-sizing;
%   		switch to using pra option until the prstper style is fixed (who knows when).
% 5/10/16: Swap suggested department/institution ordering in \affiliation lines.
% 5/12/16: Remove PACS line (deprecated).
% 6/21/16:	Switch to noshowpacs and add superscriptaddress in documentclass options.
% 
% This is my best effort to follow the Physical Review style guide plus specific changes 
% required for PERC (mostly, omitting article titles in references). This version compiles 
% with pdfLaTeX alone; using a proper .bib file changes the bibliography part at the end 
% and would require running BibTex as well.
%
% Finally, there's extra padding around section headers if you compile the bare template. 
% I believe that's because LaTeX stretches its white space to keep the floats (tables and 
% figures) near their input locations. The excess spacing goes away when a normal amount 
% of body text is filled in.

% Big list of reference documentation is here: http://journals.aps.org/revtex, 
%  see especially the APS Author Guide PDF link.

% Notes on the revtex4-2 documentclass options used: 
%  reprint does the double-column, publication ready appearance
%  prstper style formats reference numbers correctly in revtex4-2
%  titlepage creates a title page with
%  longbibliography option to show article titles (and see note in bibliography)

\documentclass[english,aps,prstper,reprint,showpacs,titlepage,longbibliography]{revtex4-2}   

\usepackage[T1]{fontenc}	% should generally be included for better accented-word behavior
\usepackage[latin9]{inputenc}	% should generally be included for better accent behavior
\usepackage{geometry}		% for controlling page margins
\geometry{verbose,tmargin=1in,bmargin=1in,lmargin=0.75in,rmargin=0.75in}	% define margins
\usepackage{graphicx}
\usepackage[above,below]{placeins}	% allows use of \FloatBarrier command to force section barriers
\usepackage{times}


% Next six lines are optional, use the hyperref package to make URLs and reference links live.
\usepackage{hyperref}  
\hypersetup{colorlinks=true,urlcolor=blue,citecolor=blue,linkcolor=blue}   
\urlstyle{same}
\pagestyle{empty}		% page numbers added later, when compiling the whole proceedings
\usepackage{enumitem}          % package for handling list formatting
\setlist{nosep}                 % Tightest spacing for lists. `noitemsep` is more relaxed

% You may use additional packages if you wish, as all you submit to the proceedings is the
% final PDF

\begin{document}

\begin{titlepage}

  \title{Title goes here, only the first word is capitalized}

  \author{First Author}
  \affiliation{Department, Some Institute, Some street address, city, state/province,
    country/postal-code} 
  \author{Second Author}
  \author{Third Author}
  \author{Fourth Author}
  \affiliation{Department, Some other Institute, Some street address, city, state/province,
    country/postal-code} 

  % \keywords{}

  \begin{abstract}
    This template is designed to provide authors of Physics Education Research Conference
    (PERC) Proceedings papers with the basis for their submission.  The format for the 2019
    Proceedings follows the PRPER style.  There are additional comments in the source code
    specific to the \LaTeX version of the template. \clearpage
  \end{abstract}
  %% Adding the `\clearpage` is the hack to make the title page.  In 2020, the proceedings is
  %% going to be double blind.  This change makes it so that we can programmatically remove the
  %% title page.  In the future, other blinding measures should be taken as well (for example,
  %% removing self-citations).  This is not needed in 2019.

  \maketitle
\end{titlepage}


\section{Introduction}
The introduction should clearly explain the nature and scope of the research as well as provide
the reader with a roadmap of the rest of the paper.

PERC proceedings submissions are limited to 4 pages.  However, the title page does not count
against that amount.  Furthermore, the references do not count against that limit either.
	
Times New Roman 10 point font is used throughout the paper with single line spacing and the
text is fully justified.  Paragraph style: indent first line three spaces (0.19''); there is no
extra line spacing between paragraphs.  This is the standard font and layout for the individual
paragraphs. % REVTeX with the options above handles this automatically.

The paper size is 8.5 $\times$ 11 inches.  Margins are 1'' top and bottom, 0.75'' left and
right.  Each column is 3.35'' with 0.3'' spacing between them.

Please do not adjust the spacing of the headers.  They should have a 12pt space before and a
6pt space after each heading. (Note: \LaTeX sets the white space itself, and sometimes pads
section headers. This is most likely to happen in documents that are underfilled with
text---like this template---and most of those problems resolve themselves as the paper reaches
a normal length.)

Information about the style of paper and answers to other submission questions can be found on
the PER--Central website~\cite{per-cent}. % Official style for references is to put them inside any following punctuation, with a space between the word and the bracketed reference number. Using a ~ as in the example here inserts a non-line-breaking space.


\subsection{Lists of all types}
You can certainly utilize lists of all types.  You can use the standard \LaTeX\ environments
for creating unordered lists
\begin{itemize}
  \item Thing
  \item Thing
  \item Thing
\end{itemize}
ordered lists
\begin{enumerate}
  \item Who
  \item What
  \item I don't know
\end{enumerate}
or even descriptions
\begin{description}
  \item[Costello] Well then who's on first?
  \item[Abbott] Yes.
  \item[Costello] I mean the fellow's name.
  \item[Abbott] Who.
  \item[Costello] The guy on first.
\end{description}

If you wish to condense the spacing between the items, you can use the \texttt{enumitem}
package and include the option \texttt{nosep}.

\section{Tables and Figures}

Figures, tables, and equations must be inserted in the text and may not be grouped at the end
of the paper. Important: miscounting of figures, tables, or equations may result from
revisions. Please double check the numbering of these elements before you submit your paper to
the volume editor.  % LaTeX handles this automatically via the \label and \ref commands.

\subsection{Figures}

If your manuscript contains figures, they are typically placed in one column (see Fig.\
\ref{fig1}).  A large figure might span two columns (see Fig.\ \ref{fig2}).  Cite all figures
in the text consecutively. The word ``Figure'' should be spelled out if it is the first word of
the sentence and abbreviated as ``Fig.'' elsewhere in the text. Place the figures as close as
possible to their first mention in the text at the top or bottom of the page.

\begin{figure}
  \includegraphics[width=0.8\linewidth]{figure1}
  \caption{Figure captions appear below the figure while table captions appear above the
    table.\label{fig1}}
\end{figure}


\subsubsection{Color figures}
Color figures are allowed since the PERC Proceedings is available only in an online format.
However, authors are encouraged to check that the figures print in black and white without loss
of clarity.

% Note: Revtex puts two-column figures at the top of a page. If you really need to change this,
% gird yourself for battle and see http://tinyurl.com/j2wnoxb.
\begin{figure*}
  \includegraphics[width=0.45\textwidth]{figure2}
  \caption{A 2-column figure. Center figure captions if they run one line only, and justify
    captions if they are multi-line.\label{fig2}}
\end{figure*}


\subsection{Tables}
Give each table a concise title and define any abbreviations in footnotes to the table. Place
tables at the top or bottom of the page, close to where they are first cited (like Table
\ref{tab1}). Center the table in the width of the text area, and center the caption as well.

Tables should not break across the page or column, but rather should be together cohesively.
In some cases, that may mean that a small amount of white space appears near the bottom of a
column.

% REVTeX provides a ruledtabular environment that automatically sizes tables, adds double
% lines, and adjusts column spacing. Use it around the regular tabular environment.  For a
% two-column table, use \begin{table*} instead (with matching \end{table*}.
\begin{table}[htbp]
  \caption{This is a sample table.\label{tab1}}
  \begin{ruledtabular}
    \begin{tabular}{cl}
      & \textbf{MS Word} \\ 
      & \textbf{US Letter Size Paper} \\ 
      \hline
      Primary Heading & Roman numerals, all capital letters \\
      First subheading & Capital letters, only first \\
      & letter of first word capitalized \\
      Second subheading & Numbers, italicized 
    \end{tabular}
  \end{ruledtabular}
\end{table}

% Formatting tweak if needed--FloatBarrier forces floats to show here, before next section
% \FloatBarrier

\section{Other issues}
\subsection{Copyright transfer}
All accepted papers will be made available under a Creative Commons Attribution 4.0 License
(CC-BY).

\subsection{Citations}
When you are preparing references, please use the \textit{Phys.\ Rev.\ PER} Style
Guide~\cite{style-guide}.  This document contains information about how references should be
formatted, including proper journal abbreviations.  Examples directly from that style guide of
a journal, book, proceedings paper, and dissertation are also provided
here~\cite{DeBoerHistory,RIOT}.  Note: These references are \textit{fake}.
In your paper, you should hyperlink to your citations as is expected by \textit{Phys.\ Rev.\
  PER}.  We do not link to non-existent papers in this template.  \LaTeX\ users, Bib\TeX\ is
your friend.

Note that for the first time in 2019, we allow references to go beyond the fourth page.  As we
no longer constrain the length of the references, you should include titles in your citations
as you typically would in \textit{Phys.\ Rev.\ PER}.

\subsection{Permissions}
To use previously published material from a book or journal, you must obtain written permission
from the owner of the rights to the material (the original publisher and/or author). It is your
responsibility to obtain permission to use copyrighted material. The executed permissions need
to be sent along with the manuscript to your volume editor. Most publishers offer submission of
permission requests online or via email, which may be the fastest and most convenient way of
receiving a reply. Some examples with relevant links are:

\begin{itemize}
  \item \url{http://www.elsevier.com/locate/permissions}
  \item \url{http://www.ieee.org}
  \item \url{http://www.nature.com}
  \item \url{http://www.sciencemag.org}
\end{itemize}

You may also use the Permission Request Form to request permission to reprint text, tables or
figures. You may complete this form and fax it to the publisher or author of the material you
wish to use. A blank form is available for download click on Forms. When the signed permission
is returned to you, please insert any necessary credit lines in your figure or table legends.


\section{Conclusions}
This template was newly updated for the 2019 PERC Proceedings.  The editors apologize if any
errors exist, and encourage you to contact them with changes and other suggestions.

\acknowledgments{Put references below the acknowledgements (and appendixes, if any).}

% \bibliographystyle{apsrev} % supercedes the longbibliography option, so leave commented out if you want to display article titles
\bibliography{bibfile} % don't include the .bib suffix

\end{document}
